% GENERADO por generar.py desde tesis.yaml. No editar a mano.
\begin{landscape}
{\small\setlength{\tabcolsep}{4pt}
\begin{longtable}{L{2.4cm}L{3.4cm}L{3.4cm}L{2.4cm}L{2.8cm}L{3.3cm}L{1.5cm}L{2.0cm}}
\caption{Matriz de operacionalización de las variables}\label{tab:operacionalizacion}\\
\toprule
\textbf{Variable} & \textbf{Definición conceptual} & \textbf{Definición operacional} & \textbf{Dimensión} & \textbf{Indicador} & \textbf{Fórmula} & \textbf{Escala} & \textbf{Instrumento} \\
\midrule
\endfirsthead
\multicolumn{8}{@{}l}{\itshape (continuación)}\\
\toprule
\textbf{Variable} & \textbf{Definición conceptual} & \textbf{Definición operacional} & \textbf{Dimensión} & \textbf{Indicador} & \textbf{Fórmula} & \textbf{Escala} & \textbf{Instrumento} \\
\midrule
\endhead
\bottomrule
\endlastfoot
\textbf{VI:} Sistema web con chatbot de WhatsApp e inteligencia artificial & EJEMPLO: aplicación web que integra la API de WhatsApp con un modelo de lenguaje para responder consultas de clientes en lenguaje natural (Autor, año). & Se implementará en la organización y se aplicará como tratamiento al grupo experimental durante el periodo de post-prueba. & --- & --- & --- & --- & --- \\
\midrule
\textbf{VD:} Gestión de atención al cliente & EJEMPLO: conjunto de actividades con las que una organización recibe, atiende y resuelve las consultas de sus clientes (Autor, año). & Se medirá mediante el tiempo de primera respuesta y la tasa de resolución, registrados en fichas de registro antes y después de implementar el sistema. & Oportunidad de la atención & Tiempo promedio de primera respuesta (TPR), en minutos & $\displaystyle TPR = \frac{\sum_{i=1}^{n} t_i}{n}$ & De razón & Ficha de registro \\
 &  &  & Eficacia de la atención & Tasa de consultas resueltas sin intervención humana (TRA), en porcentaje & $\displaystyle TRA = \frac{CR_{auto}}{CT} \times 100$ & De razón & Ficha de registro \\
\end{longtable}}
\vspace{-10pt}
\nota{Elaboración propia.}
\end{landscape}
